\documentclass{article}
\usepackage{amsmath}
\usepackage{amssymb}
\usepackage{amsthm}
\usepackage{MnSymbol}
\usepackage{xcolor}
\usepackage{parskip}
\usepackage{tikz}
\usepackage{bbm}
\usepackage{hyperref}
\usetikzlibrary{trees}
\usetikzlibrary{graphs}


\theoremstyle{definition}
\newtheorem{definition}{Definition}
\newtheorem{theorem}{Theorem}
\newtheorem{lemma}{Lemma}


\newcommand{\abs}[1]{\left| #1 \right|}


\title{lecture}
\date{\today}
\begin{document}
\maketitle
positive result:

assume realizability, assume \(m\) combinations of 3 points instead of 4 point comparisons

\(O(\sqrt{m} )\) suffices

create dependency graph, draw edges for constraint (abc), draw on (ab),(ac)

the answer will depend on \emph{arboricity}

\begin{definition}[arboricity]
    the arboricity of a graph is the min \# of forests in which edges can be partitioned
\end{definition}

if r is the arboricity, and \(S\) is the set of subgraphs of \(G\),
\begin{equation}
    r = \max_S{\frac{|E_S|}{|V_S|-1}}
\end{equation}

if a graph \(G\) has arb \(r\), then it is \(2r\) colorable

\(r\leq 2\sqrt{|E_G|}\)


\(\exists\) an ordering of \(V\in G\), \(x_1,\dots,x_n\), every \(|N(x_i)\leq 2r-1\) for \(i<j\) for \(x_j\) where \(N\) gives the neighbors of \(x\)









\end{document}